\section{if函数的基本使用}

IF 函数是 Excel 中最常用的函数之一，它可以对值和期待值进行逻辑比较。

使用逻辑函数 IF 函数时，如果条件为真，该函数将返回一个值；如果条件为假，函数将返回另一个值。


IF函数有三个参数，语法如下：

=IF(logical\_test, value\_if\_true, [value\_if\_false])

\cred{=IF(条件判断, 结果为真返回值, 结果为假返回值)}

第一参数是条件判断，比如说A1=``张三''或``A2>500''这种，结果返回TRUE或FALSE。如果判断返回TRUE那么IF函数返回值是第二参数，否则返回第三参数。

\begin{itemize}
    \item 第一参数在某些情况下可以不用比较判断符号来判断。Excel中将0认为是FALSE，非0值等价于TRUE。
    
    \cred{=IF(A1,``真'',``假'')} 
    
    如果A1为非零值，返回``真''，如果A1为0，返回``假''。
    \item 省略了第三参数，则返回第三参数值时会返回FALSE。
    
    \cred{=IF(1>2,``真'')}

    返回FALSE
    \item 逗号后面不输入参数默认返回值为0。
    
    \cred{=IF(1>2,``真'',)}

    返回0
    \item 如果要返回空白，引号内留空。
    
    \cred{=IF(1>2,``真'',``'')}
    
    \item 如果要在公式中使用文本，需要将文字用引号括起来（例如``Text''）。 唯一的例外是使用 TRUE 和 FALSE 时，Excel 能自动理解它们。
\end{itemize}

\section{多条件判断}

\begin{itemize}
    \item 多条件同时满足
    
    \cred{=IF(AND(A2>29,B2=``A''),``优秀'',``'')}

    这里用AND()函数表达了多个判断条件，只有所有判断都为TRUE才返回``优秀''。
    \item 多条件满足其一
    
    \cred{=IF(OR(A2=``A'',B2>100),``合格'',``不合格'')}

    用OR()函数来表达满足条件之一，也就是OR()函数内的判断只要有一个返回TRUE，OR()函数整体就返回TRUE。
\end{itemize}

\section{if函数的嵌套}

\cred{=IF(D2>89,"A",IF(D2>79,"B",IF(D2>69,"C",IF(D2>59,"D","F"))))}

如果单元格 D2大于 89，则学生获得 A

如果单元格 D2大于 79，则学生获得 B

如果单元格 D2大于 69，则学生获得 C

如果单元格 D2大于 59，则学生获得 D

否则，学生获得 F

\begin{itemize}
    \item 为了使长公式更易于阅读，可在编辑栏中插入换行符。 只需在将文本换到新行前按 Alt+Enter。
\end{itemize}

\section{ifs函数（Microsoft 365、Excel 2016 及更高版本）}

IFS 函数检查是否满足一个或多个条件，且返回符合第一个 TRUE 条件的值。 IFS 可以取代多个嵌套 IF 语句，并且有多个条件时更方便阅读。

通常情况下，IFS 函数的语法如下：

=IFS(logical\_test1,value\_if\_true1,logical\_test2,value\_if\_true2,logical\_test3,value\_if\_true3)

\cred{=IFS(A2>89,"A",A2>79,"B",A2>69,"C",A2>59,"D",TRUE,"F")}

如果（A2 大于 89，则返回``A''，如果 A2 大于 79，则返回``B''并以此类推，对于所有小于59的值，返回``F''）

\begin{itemize}
    \item 若要指定默认结果，请对最后一个 logical\_test 参数输入 TRUE。 如果不满足其他任何条件，则将返回相应值。
    \item 如果提供了 logical\_test 参数，但未提供相应的 value\_if\_true，则此函数显示``你为此函数输入的参数过少''错误消息。
    \item 如果 logical\_test 参数经计算解析为 TRUE 或 FALSE 以外的值，则此函数返回 \#VALUE! 错误。
    \item 如果找不到 TRUE 条件，则此函数返回 \#N/A! 错误。
\end{itemize}

\section{countif函数}

计算某个区域中满足给定条件的单元格数目

=COUNTIF(要检查哪些区域? 要查找哪些内容?)

例如：

=COUNTIF(A2:A5,``London'')

=COUNTIF(A2:A5,A4)

\begin{itemize}
    \item 任何文本条件或任何含有逻辑或数学符号的条件都必须使用双引号 (`` \quad '') 括起来。 如果条件为数字或单元格引用，则无需使用双引号。
    \item criteria 不区分大小写。
    \item 可以在 criteria 中使用通配符，即问号 (?) 和星号 (*)。 
\end{itemize}

\section{countifs函数（Microsoft 365、Excel 2016 及更高版本）}

计算某个区域中满足多个条件的单元格数目

语法

COUNTIFS(criteria\_range1, criteria1, [criteria\_range2, criteria2],…)

\begin{itemize}
    \item criteria\_range1 必需。在其中计算关联条件的第一个区域。
    \item criteria1 必需。条件的形式为数字、表达式、单元格引用或文本，它定义了要计数的单元格范围。
    
    例如，条件可以表示为 32、">32"、B4、"apples"或 "32"。
    \item criteria\_range2, criteria2, ...    可选。 附加的区域及其关联条件。
    \item 每一个附加的区域都必须与参数 criteria\_range1 具有相同的行数和列数。 这些区域无需彼此相邻。
    \item 如果条件参数是对空单元格的引用，COUNTIFS 会将该单元格的值视为 0。
    \item 可以在条件中使用通配符， 即问号 (?) 和星号 (*)。 如果要查找实际的问号或星号，请在字符前键入波形符 (\textasciitilde)。
\end{itemize}

\section{SUMIF函数}

可以使用 SUMIF 函数对范围中符合指定条件的值求和。

语法

SUMIF(range, criteria, [sum\_range])

\begin{itemize}
    \item range 范围，必需。希望通过标准评估的单元格范围。 
    \item criteria   必需。criteria 以数字、表达式、单元格参考、文本或函数的形式来定义将添加哪些单元格。
    
    例如，criteria可以表示为 32、“>32”、B5、“3？”、“苹果*”或 今天（）。
    \item sum\_range   可选。 要求和的单元格区域，如果要在range参数指定以外的其他单元格求和。sum\_range 的大小和形状应该与range相同。
\end{itemize}

例如，如果某列中含有数字，你需要对大于 5 的数值求和。 可使用以下公式：

=SUMIF(B2:B25,">5")

如果需要，可将条件应用于一个区域并对其他区域中的对应值求和。

例如，公式 =SUMIF(B2:B5, "John", C2:C5) 只对区域 C2:C5 中在区域 B2:B5 中所对应的单元格等于“John”的值求和。

\section{SUMIFS函数}

SUMIFS 函数用于计算满足多个条件的全部参数的总量。 

语法
SUMIFS(sum\_range, criteria\_range1, criteria1, [criteria\_range2, criteria2], ...)

=SUMIFS(A2:A9,B2:B9,"<>香蕉",C2:C9,"卢宁")

\begin{itemize}
    \item Sum\_range   （必需）要求和的单元格区域。
    \item Criteria\_range1   （必需）使用 Criteria1 测试的区域。
    \item Criteria1   （必需）定义将计算 Criteria\_range1 中的哪些单元格的和的条件。
    \item Criteria\_range2, criteria2, …    （可选）附加的区域及其关联条件。 
\end{itemize}